\documentclass[11pt,a4paper]{article}
\usepackage[margin=2.4cm]{geometry}
\usepackage{amsmath,amssymb,amsthm,mathtools}
\usepackage{booktabs}
\usepackage{enumitem}
\usepackage{xcolor}
\usepackage{hyperref}
\usepackage{fontspec}
\usepackage{polyglossia}
\setmainlanguage{russian}
\setotherlanguage{english}
\setmainfont{Times New Roman}
\setmonofont{Consolas}[Scale=0.90]

\hypersetup{
  colorlinks=true,
  linkcolor=blue!70!black,
  citecolor=blue!70!black,
  urlcolor=blue!70!black,
  pdftitle={Ротаторное исчисление на кубооктаэдре, Часть I: Дискретная экспонента Эйлера на Z_512, универсальный константный вронскиан и замкнутые квадратуры на кубооктаэдре O_2(3) в Lean 4},
  pdfauthor={Каратаев Дмитрий Александрович}
}

\theoremstyle{plain}
\newtheorem{theorem}{Теорема}[section]
\newtheorem{lemma}[theorem]{Лемма}
\newtheorem{proposition}[theorem]{Предложение}
\newtheorem{corollary}[theorem]{Следствие}
\theoremstyle{definition}
\newtheorem{definition}[theorem]{Определение}
\newtheorem{remark}[theorem]{Замечание}

\title{\textbf{Ротаторное исчисление на кубооктаэдре, Часть~I:\\ Дискретная экспонента Эйлера на $\mathbb{Z}_{512}$, универсальный константный вронскиан и замкнутые квадратуры на кубооктаэдре $O_2(3)$ в Lean~4}}
\author{Д.\,А.~Каратаев}
\date{}

\begin{document}
\maketitle

\begin{abstract}
В первой статье серии исследуется дискретное тригонометрическое и интегральное исчисление на трёхмерной оболочке кубооктаэдра $O_2(d)\subset\{-1,0,+1\}^d$ ($d=3$, $|V_{O_2}|=12$ вершин, $|\square|=6$ квадратных и $|\triangle|=8$ треугольных граней) и связанном с ней конечном фазовом кольце $\mathbb{Z}_{N_{\mathrm{ring}}}$ ($N_{\mathrm{ring}} = 4\,Q_\square = 512$, где $Q_\square = 2^{|\triangle|-1} = 128$ --- четверть-период). В продолжение кусочно-линейного анализа Томпсона (\texttt{arXiv:1205.3814}), где оставалась открытой задача построения интегрального исчисления без ветвления по квадрантам, показано, что пара кусочно-линейных проекций $\mathcal{E}_d(\varphi) = (\mathcal{C}_d(\varphi),\mathcal{S}_d(\varphi))$ образует точную дискретную экспоненту Эйлера на $512$-угольнике, порождающую фундаментальную пятёрку $\{0,1,i,\pi,e\}$: сдвиг фазы на четверть периода $+Q_\square(d)=+128$ действует во всех $512$ точках кольца как умножение на мнимую единицу $i$, сдвиг на полупериод $+\Phi_{\mathrm{half}}(d)=+256$ тождественно удовлетворяет соотношению Эйлера $\mathcal{E}_d(\varphi+\Phi_{\mathrm{half}}(d))+\mathcal{E}_d(\varphi)=(0,0)$ ($e^{i\pi}+1=0$), а одношаговый разностный вронскиан $\mathcal{W}_d(\varphi)=\mathcal{C}_d(\varphi)\mathcal{S}_d(\varphi+1)-\mathcal{S}_d(\varphi)\mathcal{C}_d(\varphi+1)=\langle \mathcal{E}_d(\varphi+Q_\square(d)),\Delta\mathcal{E}_d(\varphi)\rangle$ тождественно равен константе $Q_\square(d)=128$ во всех $512$ узлах кольца, включая угловые вершины $L^1$-ромба.

Доказано, что двухтактовая симплектическая рекуррентность второго порядка $\varphi_{t+1}-2\varphi_t+\varphi_{t-1}=\Phi$ объединяет элементарный круговой базис Лиувилля ($\Phi=0$) и неэлементарный квадратично-фазовый базис Френеля--Эйри ($\Phi\ne 0$), а двухкольцевая связка при собственном геометрическом модуле кубооктаэдра $k^2 = 1/\|e\|_2^2 = 1/(d-1) = 1/2$ порождает лемнискатическую триаду эллиптических функций Якоби $(\mathrm{sn}_d,\mathrm{cn}_d,\mathrm{dn}_d)$. Это позволяет вычислить в точных целых числах дискретные интегралы дифракции Френеля ($\mathcal{C}_{\mathrm{Fresnel}}=\mathcal{S}_{\mathrm{Fresnel}}=2048$), полные и неполные эллиптические интегралы II~рода, интегральный синус Дирихле ($\mathrm{Si}_3(128)=128$), интеграл Парсеваля ($2\,796\,288$) и интеграл Бесселя нулевого порядка ($\mathcal{J}_{0,3}(0)=65\,536$, $\mathcal{J}_{0,3}(1)=32\,768$ и точный корень $\mathcal{J}_{0,3}(4)=0$). Наконец, проекция $12$ вершин кубооктаэдра $O_2(3)$ на координатную ось порождает равновеликое разбиение Архимеда $(4,4,4)$, которое в сумме с центральным узлом веса $12$ автоматически выводит параболические квадратурные веса $(4,16,4)=4\times(1,4,1)$ и точный трёхмерный континуальный интеграл Римана $\iiint_{[-1,1]^3} F\,\mathrm{d}V$. Все определения и теоремы (\texttt{thm1}, \texttt{thm2}, \texttt{thm3}) машинно проверены в \textbf{Lean~4} (\texttt{0 sorryAx}, \texttt{0 Classical.choice}).

\medskip\noindent
\textbf{Ключевые слова:} кусочно-линейная тригонометрия, тождество Эйлера, разностный вронскиан, подстановка Вейерштрасса, кубооктаэдр, интегралы Френеля, лемнискатические функции Якоби, формальная верификация, Lean~4.
\end{abstract}

\section{Дискретная экспонента Эйлера, тождество \texorpdfstring{$e^{i\pi}+1=0$}{e\^{}(i pi)+1=0}, константный вронскиан и подъём Вейерштрасса на 512-угольник}
\label{sec:wronskian-euler}

На трёхмерной оболочке кубооктаэдра Архимеда $O_2(d) = \{e=(e_x,e_y,e_z)\in\{-1,0,+1\}^3 \mid e_x^2+e_y^2+e_z^2 = d-1\}$ ($d=3$, $|{\triangle}(d)|=8$, $|\square(d)|=6$, $|V_{O_2}(d)|=12$, $d_{\mathrm{in}}(\square)=4$, $d_{\mathrm{out}}(E)(d)=5$, $\|e\|_2^2=2$) и фазовом кольце $\mathbb{Z}_{N_{\mathrm{ring}}(d)}$ ($Q_\square(d)=128$, $\Phi_{\mathrm{half}}(d)=256$, $N_{\mathrm{ring}}(d)=512$, $K_P(d)=64$) определим симметричную свёртку и проекции без ветвлений:
\begin{equation}
  \label{eq:proj-cs}
  \mathrm{wrap}_d(x) := \Phi_{\mathrm{half}}(d) - \bigl((\Phi_{\mathrm{half}}(d) - x) \bmod N_{\mathrm{ring}}(d)\bigr),\quad
  \mathcal{C}_d(\varphi) := Q_\square(d) - |\mathrm{wrap}_d(\varphi)|,\quad
  \mathcal{S}_d(\varphi) := \mathcal{C}_d\bigl(\varphi - Q_\square(d)\bigr),
\end{equation}
дискретную экспоненту Эйлера $\mathcal{E}_d(\varphi) := (\mathcal{C}_d(\varphi),\mathcal{S}_d(\varphi))\in\mathbb{Z}^2$, одношаговый разностный вронскиан $\mathcal{W}_d(\varphi) := \mathcal{C}_d(\varphi)\mathcal{S}_d(\varphi+1) - \mathcal{S}_d(\varphi)\mathcal{C}_d(\varphi+1)$ и подъём Брахмагупты--Вейерштрасса $X_d(\varphi) := \mathcal{C}_d(\varphi)^2 - \mathcal{S}_d(\varphi)^2$, $Y_d(\varphi) := \|e\|_2^2\mathcal{C}_d(\varphi)\mathcal{S}_d(\varphi)$, $R_d(\varphi) := \mathcal{C}_d(\varphi)^2 + \mathcal{S}_d(\varphi)^2$. Всюду в настоящей серии статей термин \emph{носитель} (\emph{carrier}) употребляется строго в смысле классической теории моделей А.\,И.~Мальцева (\emph{«Алгебраические системы»}~\cite{Malcev1970}) как основное множество сорта в двусортной алгебраической системе $\langle\mathcal{M},\mathcal{T};\,\pi_{M\to T},\,\iota_{T\to M}\rangle$, где \textbf{дискретный $M$-носитель} $\mathcal{M}=\mathbb{Z}_{512}$ есть конечное кольцо фаз на решётке кубооктаэдра с $L^1$-проекциями $(\mathcal{C}_3,\mathcal{S}_3)$, а \textbf{непрерывный $T$-носитель} $\mathcal{T}$ --- континуальный сорт евклидовой окружности $X^2+Y^2=R^2$, поля $\mathbb{Q}\subset\mathbb{R}$ и интегралов Римана.

\begin{theorem}[Дискретное тождество Эйлера, универсальный константный вронскиан, подъём Вейерштрасса и формула Ньютона--Лейбница]
\label{thm:main-1}
При $d=3$ во всех $N_{\mathrm{ring}}(d)=512$ точках $\varphi\in\mathbb{Z}_{N_{\mathrm{ring}}(d)}$ выполнены точные тождества (\texttt{thm1\_constant\_wronskian\_weierstrass\_lift\_and\_newton\_leibniz}):
\begin{enumerate}[label=\textbf{(\roman*)},leftmargin=*]
\item $|\mathcal{C}_d(\varphi)| + |\mathcal{S}_d(\varphi)| \equiv Q_\square(d) = 128$, $\mathcal{E}_d(\varphi + Q_\square(d)) = (-\mathcal{S}_d(\varphi), +\mathcal{C}_d(\varphi)) = i\cdot\mathcal{E}_d(\varphi)$, $\mathcal{E}_d(\varphi + \Phi_{\mathrm{half}}(d)) + \mathcal{E}_d(\varphi) = (0,0)$, и во всех $512$ узлах (включая вершины излома)
  \[
    \bigl\langle \mathcal{E}_d\bigl(\varphi+Q_\square(d)\bigr),\,\Delta\mathcal{E}_d(\varphi)\bigr\rangle \;=\; \mathcal{W}_d(\varphi) \;\equiv\; Q_\square(d) \;=\; 128;
  \]
\item $X_d(\varphi)^2 + Y_d(\varphi)^2 = R_d(\varphi)^2$, $X_d(\varphi)Y_d(\varphi+1) - Y_d(\varphi)X_d(\varphi+1) = \|e\|_2^2 Q_\square(d)\mathcal{I}_d(\varphi)$ и $X_d(\varphi)X_d(\varphi+1) + Y_d(\varphi)Y_d(\varphi+1) = \mathcal{I}_d(\varphi)^2 - Q_\square(d)^2$ при $\mathcal{I}_d(\varphi) = \mathcal{C}_d(\varphi)\mathcal{C}_d(\varphi+1) + \mathcal{S}_d(\varphi)\mathcal{S}_d(\varphi+1)$;
\item первообразные $\widehat{\Phi}_{\mathcal{C}}(d,\Phi) = \sum_{\varphi=0}^{\Phi-1}\mathcal{C}_d(\varphi)$ и $\widehat{\Phi}_{\mathcal{S}}(d,\Phi) = \sum_{\varphi=0}^{\Phi-1}\mathcal{S}_d(\varphi)$ удовлетворяют $\Delta\widehat{\Phi}_{\mathcal{C}} = \mathcal{C}_d$, $\Delta\widehat{\Phi}_{\mathcal{S}} = \mathcal{S}_d$ и дают точную площадь полуволны $\widehat{\Phi}_{\mathcal{S}}(3,256) = 128^2 = 16\,384$.
\end{enumerate}
\end{theorem}

\section{Замкнутое вычисление интегралов Френеля, Якоби, Дирихле, Парсеваля и Бесселя и геометрия поясов Архимеда кубооктаэдра}
\label{sec:integrals-archimedes}

\begin{theorem}[Замкнутые значения интегралов Френеля, Якоби, Дирихле, Парсеваля и Бесселя]
\label{thm:main-2}
При $d=3$ на кольце $\mathbb{Z}_{512}$ выполнены точные интегральные соотношения (\texttt{thm2\_closed\_form\_fresnel\_jacobi\_dirichlet\_parseval\_bessel\_integrals}):
\begin{enumerate}[label=\textbf{(\roman*)},leftmargin=*]
\item для лемнискатической триады Якоби ($k^2 = 1/2$) $|\mathrm{sn}_3(\varphi)|+|\mathrm{cn}_3(\varphi)|=128$, $2\,\mathrm{dn}_3(\varphi)+|\mathrm{sn}_3(\varphi)|=256$, полный эллиптический интеграл $\mathcal{I}_{\mathrm{ell}}(3;2,3) = 163\,840$, неполный по квадранту --- $16\,384$;
\item полные интегралы дифракции Френеля совпадают и равны $\mathcal{C}_{\mathrm{Fresnel}}(3) = \mathcal{S}_{\mathrm{Fresnel}}(3) = 32\times 64 = 2048$ (выживание $32$ точек оболочки $\varphi = 8(2k+1)$ на биссектрисе $\varphi^2\equiv 64\pmod{512}$);
\item интегральный синус Дирихле $\mathrm{Si}_3(128) = \mathrm{Si}_3(256) = 128$, интеграл Парсеваля $\mathcal{I}_{\mathrm{Parseval}}(3) = \frac{512\times 32\,769}{6} = 2\,796\,288$;
\item дискретный интеграл Бесселя $\mathcal{J}_{0,3}(0) = 65\,536$, $\mathcal{J}_{0,3}(1) = 32\,768$, и точный корень $\mathcal{J}_{0,3}(4) = 0$.
\end{enumerate}
\end{theorem}

\begin{theorem}[Геометрический вывод континуального интеграла Римана, весов Симпсона $(1,4,1)/6$, тождества куба вронскиана для $\pi(n)$ и 3D-ротатора Эйлера из кубооктаэдра]
\label{thm:main-3}
При $d=3$ на оболочке кубооктаэдра $O_2(3)$ и кольце $\mathbb{Z}_{512}$ выполнены соотношения (\texttt{thm3\_continuum\_integration\_from\_cuboctahedron\_archimedes\_layers}):
\begin{enumerate}[label=\textbf{(\roman*)},leftmargin=*]
\item проекция $12$ вершин $O_2(3)$ на любую ось $e_x\in\{-1,0,+1\}$ образует три равновеликих пояса Архимеда $(4,4,4)$, которые в сумме с центральным узлом веса $12$ в операторе $\widehat{\mathcal{I}}_{O_2(3)}(F) = 12 F(0,0,0) + \sum_{e\in O_2(3)} F(e)$ собираются в веса $(4,16,4) = 4\times(1,4,1)$ и воспроизводят континуальный интеграл Римана $\frac{1}{24}\widehat{\mathcal{I}}_{O_2(3)}(x^p) = \frac{1}{2}\int_{-1}^{+1}x^p\,\mathrm{d}x$ ($0\le p\le 3$) и $\frac{1}{24}\widehat{\mathcal{I}}_{O_2(3)}(x^2+y^2+z^2) = 1$;
\item рациональная вилка $\pi_{\mathrm{out}}(n) - \pi_{\mathrm{in}}(n) = \sum_{j=0}^{n/2-1}\frac{W_s(js)^3}{I_s(js)R_d(js)R_d((j+1)s)} > 0$ сжимается в $4$ раза при каждом удвоении $n\mapsto 2n$, а рациональный оператор синуса $\sin_{\mathrm{rat}}(3;p,q) = \frac{4p(2q-p)}{5q^2-p(2q-p)}$ на долях прямого угла $\{1/9,1/3,1/2,2/3,1\}$ возвращает $\frac{17}{97},\frac{1}{2},\frac{12}{17},\frac{32}{37},1$;
\item моменты кубооктаэдра равны $\sum_{e\in O_2(3)}(\mathbf{k}\cdot e)^2 = 8|\mathbf{k}|^2$ и $\sum_{e\in O_2(3)}(\mathbf{k}\cdot e)^4 = 12|\mathbf{k}|^4 - 4\sum_{i=1}^3 k_i^4$, а трёхмерный ротатор Эйлера $\mathcal{R}_{\mathrm{Euler}}(\varphi_1,\varphi_2,\varphi_3)$ сохраняет сферическую норму в чистых целых числах.
\end{enumerate}
\end{theorem}

\section{Формальная верификация в Lean~4}
Формальное доказательство Теорем~\ref{thm:main-1}--\ref{thm:main-3} (\texttt{thm1}, \texttt{thm2}, \texttt{thm3}) содержится в модуле \texttt{RotatorCalculusPaper.lean} в репозитории \url{https://github.com/AI-Guiders/scientific-papers} (\texttt{0 sorryAx}, \texttt{0 Classical.choice}).

\begin{thebibliography}{9}
\bibitem{Malcev1970}
А.\,И.~Мальцев, \emph{Алгебраические системы}, М.: Наука, 1970 (англ.\ пер.: A.\,I.~Mal'cev, \emph{Algebraic Systems}, Springer-Verlag, 1973).
\bibitem{Krause1986}
E.\,F.~Krause, \emph{Taxicab Geometry}, Dover Publications, New York, 1986.
\bibitem{Thompson2012}
K.\,P.~Thompson, ``Taxicab calculus: Trig derivatives,'' \texttt{arXiv:1205.3814v2} [math.MG], 2012.
\bibitem{Moura2021}
L.~de~Moura and S.~Ullrich, ``The Lean~4 theorem prover and programming language,'' in \emph{CADE~28}, LNCS vol.~12699, Springer, 2021, pp.~625--635.
\end{thebibliography}

\end{document}
